%% toc: P15 166 166.3
% Continuous entropy relative to a reference measure and a coordinate; Shannon's Section 20 properties and coordinate changes.
\chapter{Continuous Entropy, Coordinates and Reference Measures}
\label{sec:p15-w02b-p15-042-continuous-entropy}

\emph{Purpose.} Shannon's 1948 theory defines the entropy of a continuous
distribution by analogy with the discrete formula and lists nine properties
of it \cite[sec.~20, pp.~35--38]{Shannon1948MathematicalTheory}. The analogy hides a
choice. The formula integrates a density, and a density is a derivative with
respect to a reference measure; in Shannon's setting that measure is volume in
a chosen coordinate system. This chapter makes the reference measure explicit,
states when the entropy integral exists, proves each of Shannon's nine
continuous properties under stated hypotheses, and proves the coordinate-change
formula on which the coordinate dependence rests.

\emph{Inputs.} A sigma-finite measure space \((S,\Sigma,\mu)\), usually
\(\mathbb R^n\) with Borel Lebesgue measure \(\lambda^n\) of
Definition~\ref{def:p15-probability-density}; a probability law with a
density relative to it; for coordinate changes, a bijection together with a
Jacobian factor (Definition~\ref{def:p15-w02b-p15-042-coordinate-change}).
\emph{Conclusions.} Whether the entropy is finite, \(+\infty\), \(-\infty\) or
undefined; Shannon's bounds, maximizers, joint and conditional identities and
averaging inequality; the exact change of entropy under a change of reference
measure or coordinates; and which derived quantities do not change.

The discrete entropy \(-\sum p_i\log p_i\) of Shannon's Section~6 and its
characterization are treated in the separate chapter on discrete entropy in
this part; in the present notation it is the case in which \(\mu\) is
counting measure. The differential entropy used for measurement models in
Chapter~\ref{sec:p15-measurement-uncertainty} is the case \(\mu=\lambda\)
on the real line.

\section{Entropy relative to a reference measure}
\label{sec:p15-w02b-p15-042-reference}

Shannon notes that choosing the base of the logarithm amounts to choosing a
unit of information: base 2 gives bits and base \(e\) gives natural units, and
passing from base \(a\) to base \(b\) multiplies by \(\log_b a\)
\cite[pp.~1--2]{Shannon1948MathematicalTheory}. All logarithms below are natural, so
entropies are in natural units; dividing by \(\log 2\) converts them to bits.

\begin{hbdefinition}{def:p15-w02b-p15-042-reference-entropy}{Entropy relative to a reference measure}
Let \((S,\Sigma,\mu)\) be a sigma-finite measure space and write
\(\eta(t)\defeq-t\log t\) for \(0<t<\infty\), with \(\eta(0)\defeq0\).
A probability measure \(P\) on \(\Sigma\) has a \emph{\(\mu\)-density}
when there is a measurable \(p:S\to[0,\infty)\) with
\(P(A)=\int_A p\,d\mu\) for every \(A\in\Sigma\).
The \emph{entropy of \(P\) relative to \(\mu\)} is
\[
 H_\mu(P)\defeq\int_S\eta(p)\,d\mu
          =-\int_S p\log p\,d\mu ,
\]
defined exactly when at least one of \(\int\eta(p)^+\,d\mu\) and
\(\int\eta(p)^-\,d\mu\) is finite; its value then lies in
\([-\infty,+\infty]\). If \(P\) has no \(\mu\)-density, or both integrals are
infinite, \(H_\mu(P)\) is undefined. Two densities of \(P\) agree
\(\mu\)-a.e.\ by the uniqueness part of \cref{thm:p15-radon-nikodym}, so the
value does not depend on the version chosen. We also write \(H_\mu(p)\).

\emph{Shannon's case.} For \(S=\mathbb R^n\) with coordinates
\(x=(x_1,\ldots,x_n)\) and \(\mu=\lambda^n\), Shannon writes
\(H=-\int p(x)\log p(x)\,dx\), and \(H(x)\) for the entropy of the chance
variable \(x\) \cite[sec.~20, p.~35]{Shannon1948MathematicalTheory}; the argument
\(x\) labels the variable and is not a function argument
\cite[sec.~6, p.~11]{Shannon1948MathematicalTheory}.
If \(x\) and \(y\) take values in \(\mathbb R^{n_1}\) and \(\mathbb R^{n_2}\)
with joint density \(p(x,y)\) relative to \(\lambda^{n_1+n_2}\), with marginal
densities \(p(x)=\int p(x,y)\,dy\) and \(p(y)=\int p(x,y)\,dx\), then
\[
 H(x,y)=-\iint p(x,y)\log p(x,y)\,dx\,dy,\qquad
 H_x(y)=-\iint p(x,y)\log\frac{p(x,y)}{p(x)}\,dx\,dy,
\]
and \(H_y(x)\) is formed with \(p(y)\) in place of \(p(x)\)
\cite[sec.~20, p.~35]{Shannon1948MathematicalTheory}. The integrand of \(H_x(y)\) is
taken to be zero where \(p(x,y)=0\). The same formulas define these
quantities for product reference measures \(\mu_1\times\mu_2\).
\index{entropy!relative to a reference measure}\index{differential entropy}
\end{hbdefinition}

The reference measure is part of the definition. Replacing \(\mu\) by another
measure with the same null sets changes the density and, as
\cref{thm:p15-w02b-p15-042-reference-change} shows, the entropy. The
following inequality, a form of Gibbs' inequality, is the single tool behind
the bounds and maximizers of this chapter.

\begin{hblemma}{thm:p15-w02b-p15-042-gibbs}{proved=annex:p15-w02b-p15-042-gibbs}{Comparison with a subprobability density}
Let \(p,q:S\to[0,\infty)\) be measurable with \(\int p\,d\mu=1\) and
\(\int q\,d\mu\leq1\). Suppose that \(p\log q\) is \(\mu\)-integrable, with
\(0\log0=0\); in particular \(\mu(\{p>0,\ q=0\})=0\). Then \(H_\mu(p)\) is
defined, belongs to \([-\infty,\infty)\), and
\[
 H_\mu(p)\leq-\int p\log q\,d\mu ,
\]
with equality if and only if \(p=q\) \(\mu\)-a.e.
\end{hblemma}

\begin{hbtheorem}{thm:p15-w02b-p15-042-reference-change}{proved=annex:p15-w02b-p15-042-reference-change}{Changing the reference measure}
Let \(f:S\to(0,\infty)\) be measurable and let \(\nu=f\mu\), that is,
\(\nu(A)=\int_A f\,d\mu\). If \(P\) has \(\mu\)-density \(p\), then \(P\) has
\(\nu\)-density \(p/f\). If \(H_\mu(P)\) is defined and
\(\int p\,|\log f|\,d\mu<\infty\), or if \(H_\mu(P)\) is finite and
\(\int p\log f\,d\mu\) exists in \([-\infty,\infty]\), then \(H_\nu(P)\) is
defined and
\[
 H_\nu(P)=H_\mu(P)+\int p\log f\,d\mu .
\]
\end{hbtheorem}

\begin{hbtheorem}{thm:p15-w02b-p15-042-defined-cases}{proved=annex:p15-w02b-p15-042-defined-cases}{Finite, infinite and undefined entropies}
\begin{enumerate}
\item If \(\mu(S)<\infty\), then \(H_\mu(P)\) is defined for every \(P\)
      with a \(\mu\)-density and \(H_\mu(P)<\infty\).
\item The function \(p_1(x)=\log2/(x(\log x)^2)\) on \((0,\tfrac12)\), zero
      elsewhere, is a density relative to \(\lambda\) with
      \(H_\lambda(p_1)=-\infty\).
\item The function \(p_2(x)=1/(x(\log x)^2)\) on \((e,\infty)\), zero
      elsewhere, is a density relative to \(\lambda\) with
      \(H_\lambda(p_2)=+\infty\).
\item For \(p_3=\tfrac12p_1+\tfrac12p_2\), both \(\int\eta(p_3)^+\,d\lambda\)
      and \(\int\eta(p_3)^-\,d\lambda\) are infinite, so
      \(H_\lambda(p_3)\) is undefined.
\item A law with an atom, such as the Dirac law \(\delta_c\), and the
      Cantor law of \cref{thm:p15-cantor-singular-law} have no density
      relative to \(\lambda\); their entropy relative to \(\lambda\) is
      undefined.
\end{enumerate}
\end{hbtheorem}

Item~1 says that a bounded region rules out the value \(+\infty\) but not the
value \(-\infty\); items~2--4 show that no finiteness can be presumed on the
real line. Item~5 marks the boundary of Shannon's definition: a distribution
must have a density before its continuous entropy means anything.

\section{Shannon's properties of continuous entropy}
\label{sec:p15-w02b-p15-042-properties}

Shannon states that continuous entropies have most, but not all, of the
properties of the discrete case, and lists them as items 1--9
\cite[sec.~20, pp.~35--38]{Shannon1948MathematicalTheory}. Items 1--7 are stated
below with the hypotheses under which they are proved; items 8 and 9 are the
subject of Section~\ref{sec:p15-w02b-p15-042-coordinates}.

\begin{hbtheorem}{thm:p15-w02b-p15-042-volume}{proved=annex:p15-w02b-p15-042-volume}{Item 1: the bound by the logarithm of the volume}
Let \(A\in\Sigma\) with \(0<v=\mu(A)<\infty\), and let \(P\) have a
\(\mu\)-density \(p\) with \(P(A)=1\). Then \(H_\mu(P)\) is defined and
\[
 H_\mu(P)\leq\log v ,
\]
with equality if and only if \(p=1/v\) \(\mu\)-a.e.\ on \(A\) (and \(p=0\)
\(\mu\)-a.e.\ off \(A\)). The value \(-\infty\) is possible.
\end{hbtheorem}

For \(\mu=\lambda^n\) this is Shannon's item~1: when \(x\) is confined to a
volume \(v\), the maximum entropy is \(\log v\), attained by the constant
density \(1/v\) \cite[sec.~20, p.~35]{Shannon1948MathematicalTheory}.

\begin{hbtheorem}{thm:p15-w02b-p15-042-joint}{proved=annex:p15-w02b-p15-042-joint}{Item 2: subadditivity and its equality case}
Let \(\mu_1,\mu_2\) be sigma-finite and let the pair \((x,y)\) have a density
\(p(x,y)\) relative to \(\mu_1\times\mu_2\), with marginal densities
\(p(x)\), \(p(y)\) obtained by integrating out the other variable.
Suppose \(H(x)\) and \(H(y)\) are finite. Then \(H(x,y)\) is defined, belongs
to \([-\infty,\infty)\), and
\[
 H(x,y)\leq H(x)+H(y),
\]
with equality if and only if \(p(x,y)=p(x)p(y)\) for
\((\mu_1\times\mu_2)\)-almost every \((x,y)\), which holds if and only if
\(x\) and \(y\) are independent.
\end{hbtheorem}

Shannon's item~2 states the same inequality and equality condition, the
equality holding apart possibly from a set of probability zero
\cite[sec.~20, p.~36]{Shannon1948MathematicalTheory}; the finiteness of the two
marginal entropies is the hypothesis that makes the right-hand side
meaningful.

\begin{hbtheorem}{thm:p15-w02b-p15-042-averaging}{proved=annex:p15-w02b-p15-042-averaging}{Item 3: generalized averaging does not decrease entropy}
Let \(a:S\times S\to[0,\infty)\) be product measurable with
\[
 \int a(x,y)\,d\mu(x)=1\ \text{for every }y,\qquad
 \int a(x,y)\,d\mu(y)=1\ \text{for every }x .
\]
Let \(p\) be a \(\mu\)-density with \(H_\mu(p)\) finite and set
\(p'(y)=\int a(x,y)p(x)\,d\mu(x)\). Then \(p'\) is a \(\mu\)-density,
\(H_\mu(p')\) is defined, and \(H_\mu(p')\geq H_\mu(p)\); the value
\(+\infty\) is not excluded.
\end{hbtheorem}

This is Shannon's item~3 with \(\mu=\lambda^n\)
\cite[sec.~20, p.~36]{Shannon1948MathematicalTheory}. The kernel condition in both
variables is what is used: integration over \(x\) makes \(a(\cdot,y)p\) an
average of \(p\), and integration over \(y\) returns the total mass of \(p\).

\begin{hbtheorem}{thm:p15-w02b-p15-042-conditional}{proved=annex:p15-w02b-p15-042-conditional}{Item 4: joint and conditional entropies}
In the setting of \cref{thm:p15-w02b-p15-042-joint}, suppose \(H(x)\),
\(H(y)\) and \(H(x,y)\) are finite. Then \(H_x(y)\) and \(H_y(x)\) are finite,
\[
 H(x,y)=H(x)+H_x(y)=H(y)+H_y(x),
\]
and \(H_x(y)\leq H(y)\), with equality if and only if \(x\) and \(y\) are
independent. The conditional entropy \(H_x(y)\) can be negative: for \((x,y)\)
uniform on \((0,1)\times(0,\tfrac12)\) relative to \(\lambda^2\),
\(H_x(y)=-\log2\).
\end{hbtheorem}

Shannon's item~4 records the two decompositions and the inequality
\cite[sec.~20, p.~36]{Shannon1948MathematicalTheory}. He also observes that the
entropy of a continuous distribution can be negative, the scale of
measurement fixing an arbitrary zero at the uniform distribution on a unit
volume \cite[sec.~20, p.~38]{Shannon1948MathematicalTheory}; the last sentence of the
theorem is an instance for a conditional entropy.

For the Gaussian items, let \(a=(a_{ij})\) be a symmetric positive definite
\(n\times n\) matrix, write \(|a_{ij}|\) for its determinant as Shannon does,
and set
\[
 q_a(x_1,\ldots,x_n)=\frac{|a_{ij}|^{1/2}}{(2\pi)^{n/2}}
   \exp\Bigl(-\tfrac12\sum_{i,j}a_{ij}x_ix_j\Bigr).
\]
For \(n=1\) and \(a_{11}=1/\sigma^2\) this is the normal density
\(\phi_{0,\sigma^2}\) of Definition~\ref{def:p15-normal-law}.

\begin{hbtheorem}{thm:p15-w02b-p15-042-gaussian}{proved=annex:p15-w02b-p15-042-gaussian}{Items 5 and 6: Gaussian maximizers and Gaussian entropies}
\begin{enumerate}
\item Let \(c\in\mathbb R\) and \(\sigma>0\). Every density \(p\) relative to
      \(\lambda\) with \(\int(x-c)^2p(x)\,dx=\sigma^2\) has a defined
      entropy, and
      \[
       H(x)\leq\log\sqrt{2\pi e}\,\sigma ,
      \]
      with equality if and only if \(p=\phi_{c,\sigma^2}\) a.e.
      For \(c=0\) this is Shannon's constraint \(\sigma^2=\int p(x)x^2\,dx\);
      for \(c\) equal to the mean of \(p\) it fixes the standard deviation.
\item The function \(q_a\) is a probability density relative to
      \(\lambda^n\), and
      \[
       \int x_ix_jq_a(x)\,dx=A_{ij},
      \]
      where \((A_{ij})\) is the inverse matrix of \((a_{ij})\).
\item Let \((A_{ij})\) be symmetric positive definite and \((a_{ij})\) its
      inverse. Every density \(p\) relative to \(\lambda^n\) with
      \(\int x_ix_jp(x)\,dx=A_{ij}\) for all \(i,j\) has a defined entropy,
      and
      \[
       H(x)\leq\log\,(2\pi e)^{n/2}|a_{ij}|^{-1/2},
      \]
      with equality if and only if \(p=q_a\) a.e.
\item If \((A_{ij})\) is symmetric but not positive definite, no density
      relative to \(\lambda^n\) has these second moments.
\item The entropy of \(\phi_{c,\sigma^2}\) is
      \(\log\sqrt{2\pi e}\,\sigma\), and the entropy of \(q_a\) is
      \[
       \log\,(2\pi e)^{n/2}|a_{ij}|^{-1/2}.
      \]
\end{enumerate}
\end{hbtheorem}

Shannon obtains items 5 and 6 by the calculus of variations: the stationarity
condition for the constrained integral forces \(\log p\) to be a quadratic
polynomial, and adjusting constants gives the Gaussian
\cite[sec.~20, pp.~36--37]{Shannon1948MathematicalTheory}. A stationarity condition
alone does not show that the stationary density is a maximum, nor that it is
the only one. The Gibbs comparison of \cref{thm:p15-w02b-p15-042-gibbs}
supplies both, and item~4 of the theorem identifies the singular moment matrices for which
the problem has no admissible density at all.

\begin{hbtheorem}{thm:p15-w02b-p15-042-exponential}{proved=annex:p15-w02b-p15-042-exponential}{Item 7: the exponential maximizer on a half line}
Let \(a>0\). Every density \(p\) relative to \(\lambda\) with \(p=0\) a.e.\
on \((-\infty,0]\) and \(\int_0^\infty x\,p(x)\,dx=a\) has a defined entropy,
and
\[
 H(x)\leq\log ea ,
\]
with equality if and only if \(p(x)=a^{-1}e^{-x/a}\) for almost every
\(x>0\).
\end{hbtheorem}

This is Shannon's item~7 \cite[sec.~20, p.~37]{Shannon1948MathematicalTheory}.

\section{Coordinates and the change of entropy}
\label{sec:p15-w02b-p15-042-coordinates}

Shannon's item~8 states the one important difference between the two cases.
In his words, ``In the continuous case the measurement is relative to the
coordinate system'' \cite[sec.~20, p.~37]{Shannon1948MathematicalTheory}. The
entropy is a measure of randomness relative to the standard that gives equal
weight to equal volume elements of the chosen coordinates
\cite[sec.~20, p.~38]{Shannon1948MathematicalTheory}. The precise content is a
change-of-variables identity. We state it as a hypothesis on the coordinate
change and then prove it for affine maps, products and \(C^1\)
diffeomorphisms.

\begin{hbdefinition}{def:p15-w02b-p15-042-coordinate-change}{Coordinate change with a Jacobian factor}
Let \(U,V\subseteq\mathbb R^n\) be Borel sets. A \emph{coordinate change with
Jacobian factor} \(j\) is a bijection \(T:U\to V\) such that \(T\) and
\(T^{-1}\) are Borel measurable, together with a Borel function
\(j:U\to(0,\infty)\), such that
\[
 \int_V g(y)\,dy=\int_U g(T(x))\,j(x)\,dx
\]
for every Borel \(g:V\to[0,\infty]\). In Shannon's notation the Jacobian
\(J(x/y)\) of the old coordinates with respect to the new, evaluated at
\(y=T(x)\), is \(1/j(x)\); here \(J(x/y)\) denotes the absolute value of
that determinant.

The hypothesis holds in the following cases:
\begin{enumerate}
\item \(T(x)=Mx+b\) on \(U=V=\mathbb R^n\), with \(M\) an invertible real
      matrix, \(b\in\mathbb R^n\) and \(j=|\det M|\)
      (\cref{thm:p15-w02b-p15-042-linear-lebesgue});
\item \(n=1\), \(U,V\) open intervals, and \(T\) a \(C^1\) bijection with
      \(T'>0\) everywhere or \(T'<0\) everywhere, with \(j=|T'|\)
      (\cref{thm:p15-criteria-real-integration});
\item products: if \((T_1,j_1)\) on \(U_1\subseteq\mathbb R^{n_1}\) and
      \((T_2,j_2)\) on \(U_2\subseteq\mathbb R^{n_2}\) satisfy the hypothesis,
      so does \(T_1\times T_2\) on \(U_1\times U_2\) with factor
      \(j_1(x)j_2(y)\) (\cref{thm:p15-w02b-p15-042-linear-lebesgue});
\item \(U,V\subseteq\mathbb R^n\) open and \(T\) a \(C^1\) diffeomorphism,
      that is, a bijection such that \(T\) and \(T^{-1}\) both have
      continuous first partial derivatives, with \(j=|\det DT|\), where
      \(DT(x)\) is the matrix of first partial derivatives of \(T\) at \(x\)
      (\cref{thm:p15-w02b-p15-042-diffeomorphism}).
\end{enumerate}
Item~4 asks for the inverse to be \(C^1\). For a \(C^1\) bijection between
open sets whose Jacobian determinant never vanishes, the inverse function
theorem supplies that property; the inverse function theorem is not proved
in this handbook, so for such a map the continuity of the derivative of
\(T^{-1}\) is to be checked before item~4 is used.
\index{coordinate change!Jacobian factor}
\end{hbdefinition}

\begin{hblemma}{thm:p15-w02b-p15-042-linear-lebesgue}{proved=annex:p15-w02b-p15-042-linear-lebesgue}{Linear and product coordinate changes}
Let \(M\) be an invertible real \(n\times n\) matrix and \(b\in\mathbb R^n\).
For every Borel \(g:\mathbb R^n\to[0,\infty]\),
\[
 \int g(Mx+b)\,dx=\frac1{|\det M|}\int g(z)\,dz ,
\]
and consequently \(\lambda^n(MA+b)=|\det M|\,\lambda^n(A)\) for Borel \(A\).
The product statement in item~3 of
Definition~\ref{def:p15-w02b-p15-042-coordinate-change} holds.
\end{hblemma}

\begin{hblemma}{thm:p15-w02b-p15-042-diffeomorphism}{proved=annex:p15-w02b-p15-042-diffeomorphism}{Change of variables for a \(C^1\) diffeomorphism}
Let \(U,V\subseteq\mathbb R^n\) be open and let \(T:U\to V\) be a \(C^1\)
diffeomorphism. Then \(\det DT(x)\neq0\) for every \(x\in U\), and for every
Borel \(g:V\to[0,\infty]\)
\[
 \int_V g(y)\,dy=\int_U g(T(x))\,\bigl|\det DT(x)\bigr|\,dx .
\]
In particular \(\lambda^n(T(A))=\int_A|\det DT|\,dx\) for every Borel
\(A\subseteq U\), and \((T,|\det DT|)\) is a coordinate change in the sense
of Definition~\ref{def:p15-w02b-p15-042-coordinate-change}.
\end{hblemma}

\begin{hbtheorem}{thm:p15-w02b-p15-042-coordinate-change-thm}{proved=annex:p15-w02b-p15-042-coordinate-change-thm}{Item 8: entropy in new coordinates}
Let \((T,j)\) be a coordinate change from \(U\) to \(V\), let \(x\) have a
density \(p\) relative to \(\lambda^n\) with \(P(x\in U)=1\), and let
\(y=T(x)\). Then \(y\) has the density
\(p_y(t)=p(T^{-1}t)/j(T^{-1}t)\) on \(V\), zero elsewhere. If \(H(x)\) is
defined and \(\int p\,|\log j|\,dx<\infty\), or if \(H(x)\) is finite and
\(\int p\log j\,dx\) exists in \([-\infty,\infty]\), then \(H(y)\) is defined
and
\[
 H(y)=H(x)+\int p(x)\log j(x)\,dx
     =H(x)-\int p(x)\log J\Bigl(\frac{x}{y}\Bigr)\,dx .
\]
\end{hbtheorem}

The second form is Shannon's: the new entropy is the old one less the
expected logarithm of the Jacobian \cite[sec.~20, pp.~37--38]{Shannon1948MathematicalTheory}.
In the language of \cref{thm:p15-w02b-p15-042-reference-change}, passing to
\(y\) replaces the reference measure \(\lambda^n\) on the \(x\)-coordinates by
the measure \(j\,\lambda^n\), whose image under \(T\) is \(\lambda^n\) on the
\(y\)-coordinates.

\begin{hbcorollary}{thm:p15-w02b-p15-042-linear-change}{proved=annex:p15-w02b-p15-042-linear-change}{Item 9: linear coordinates}
Let \((a_{ij})\) be an invertible real \(n\times n\) matrix and
\(y_j=\sum_ia_{ij}x_i\). If \(H(x)\) is defined, then \(H(y)\) is defined and
\[
 H(y)=H(x)+\log\bigl|\det(a_{ij})\bigr| .
\]
In particular \(H(y)=H(x)\) for a rotation, a translation, or any
coordinate change whose Jacobian factor is identically \(1\).
\end{hbcorollary}

Shannon writes the Jacobian as \(|a_{ij}|^{-1}\) and the result as
\(H(y)=H(x)+\log|a_{ij}|\), with \(|a_{ij}|\) his determinant; the absolute
value of the determinant is what enters when the transformation reverses
orientation \cite[sec.~20, p.~38]{Shannon1948MathematicalTheory}.

\begin{hbcorollary}{thm:p15-w02b-p15-042-difference-invariance}{proved=annex:p15-w02b-p15-042-difference-invariance}{Differences that do not depend on coordinates}
Let \((x,y)\) have a density relative to \(\lambda^{n_1+n_2}\) with
\(H(x)\), \(H(y)\) and \(H(x,y)\) finite. Let \((T_1,j_1)\) and
\((T_2,j_2)\) be coordinate changes on \(\mathbb R^{n_1}\) and
\(\mathbb R^{n_2}\) with \(P(x\in U_1)=P(y\in U_2)=1\),
\(\int p(x)|\log j_1(x)|\,dx<\infty\) and
\(\int p(y)|\log j_2(y)|\,dy<\infty\). For \(x'=T_1(x)\), \(y'=T_2(y)\), the
entropies \(H(x')\), \(H(y')\), \(H(x',y')\) are finite and
\[
 H(x')+H(y')-H(x',y')=H(x)+H(y)-H(x,y)=H(x)-H_y(x).
\]
\end{hbcorollary}

This is the precise form of Shannon's remark that rates and capacities depend
on a difference of two entropies, each term changing by the same amount
\cite[sec.~20, p.~38]{Shannon1948MathematicalTheory}. The hypothesis matters:
coordinates that mix \(x\) with \(y\) are not of this product form, and the
difference is then not covered by the corollary.

\begin{hbtheorem}{thm:p15-w02b-p15-042-boundaries}{proved=annex:p15-w02b-p15-042-boundaries}{Coordinate-relative values, rankings and divergence}
Take \(U=V=(0,1)\) and \(y=x^2\), a coordinate change with factor
\(j(x)=2x\).
\begin{enumerate}
\item If \(x\) is uniform on \((0,1)\), then \(H(x)=0\) and
      \(H(y)=\log2-1\approx-0.307\).
\item If \(x\) has density \(2x\) on \((0,1)\), then \(y\) is uniform on
      \((0,1)\), \(H(x)=\tfrac12-\log2\approx-0.193\) and \(H(y)=0\).
\end{enumerate}
Thus the first distribution has the larger entropy in the \(x\)-coordinate
and the smaller in the \(y\)-coordinate.
\begin{enumerate}\setcounter{enumi}{2}
\item If \(x\) is uniform on \((0,1)\) and \(y=e^{-1/x}\), a coordinate change
      from \((0,1)\) onto \((0,e^{-1})\), then \(H(x)=0\) and
      \(H(y)=-\infty\).
\item If \(z_1,z_2\) are independent standard normal variables, \(|\rho|<1\),
      \(\sigma_1,\sigma_2>0\), and
      \(x_1=\sigma_1z_1\), \(x_2=\sigma_2(\rho z_1+\sqrt{1-\rho^2}\,z_2)\),
      then \(x_i\) has law \(\operatorname{normal}(0,\sigma_i^2)\),
      \(H(x_1,x_2)=\log\bigl(2\pi e\,\sigma_1\sigma_2\sqrt{1-\rho^2}\bigr)\),
      and
      \[
       H(x_1)+H(x_2)-H(x_1,x_2)=-\tfrac12\log(1-\rho^2).
      \]
\end{enumerate}
\end{hbtheorem}

Items 1 and 2 are instances of a family that covers every power coordinate on
the unit interval.

\begin{hbcorollary}{thm:p15-w02b-p15-042-power-coordinates}{proved=annex:p15-w02b-p15-042-power-coordinates}{Uniform laws in power coordinates}
Let \(\ell\) take values in \((0,1)\), let \(a,b>0\), and suppose that
\(\ell^a\) is uniform on \((0,1)\). Put \(r=b/a\). Then
\(P(\ell\leq c)=c^a\) for \(0<c<1\), the map \(t\mapsto t^r\) is a
coordinate change of \((0,1)\) onto itself with factor \(j(t)=rt^{r-1}\),
and
\[
 H(\ell^b)=\log r+1-r\leq0 ,
\]
with equality if and only if \(a=b\). Consequently, among the laws for which
some power \(\ell^a\) is uniform, the law that is uniform in \(\ell^b\) is the
only one with the largest entropy in the coordinate \(\ell^b\).
\end{hbcorollary}

\section{Procedure and worked boundary}
\label{sec:p15-w02b-p15-042-procedure}

\begin{hbmethod}{def:p15-w02b-p15-042-method}{Computing and reporting a continuous entropy}
\begin{enumerate}
\item Fix the space, the reference measure and, on \(\mathbb R^n\), the
      coordinate system in which it is volume; fix the logarithm base.
\item Establish that the law has a density relative to that measure
      (\cref{thm:p15-w02b-p15-042-defined-cases}, item~5).
\item Decide whether \(\int\eta(p)^+\) and \(\int\eta(p)^-\) are finite, and
      report the entropy as finite, \(+\infty\), \(-\infty\) or undefined.
\item Before using a property, verify its hypotheses: a finite reference
      volume for item~1, finite marginal entropies for items 2 and 4, a
      doubly normalized kernel for item~3, a positive definite moment matrix
      for item~5.
\item For new coordinates, establish the Jacobian identity of
      Definition~\ref{def:p15-w02b-p15-042-coordinate-change} (for a
      \(C^1\) diffeomorphism of open sets it holds with
      \(j=|\det DT|\) by \cref{thm:p15-w02b-p15-042-diffeomorphism}) and the
      integrability of \(p\log j\); then apply
      \cref{thm:p15-w02b-p15-042-coordinate-change-thm}.
\item Report every entropy with its reference measure, coordinates and unit.
      Compare entropies across coordinate systems only through quantities
      shown to be invariant, such as the difference in
      \cref{thm:p15-w02b-p15-042-difference-invariance}.
\end{enumerate}
\end{hbmethod}

\begin{hbexample}{def:p15-w02b-p15-042-units-example}{A change of units and a change of coordinates}
\emph{Units.} Let a displacement \(x\), in metres, have law
\(\operatorname{normal}(0,\sigma^2)\) with \(\sigma=0.002\). By
\cref{thm:p15-w02b-p15-042-gaussian}, item~5,
\[
 H(x)=\tfrac12\log(2\pi e)+\log0.002\approx1.4189-6.2146=-4.7957
\]
natural units. The same displacement in millimetres is \(y=1000x\). By
\cref{thm:p15-w02b-p15-042-linear-change},
\(H(y)=H(x)+\log1000\approx-4.7957+6.9078=2.1121\), which is also
\(\tfrac12\log(2\pi e)+\log2\), the entropy of a normal law with standard
deviation \(2\). The sign of the entropy changed with the unit; neither sign
says anything about the displacement itself.

\emph{Dependence.} Let two displacements \((x_1,x_2)\) be as in
\cref{thm:p15-w02b-p15-042-boundaries}, item~4, with \(\rho=0.6\). Then
\(H(x_1)+H(x_2)-H(x_1,x_2)=-\tfrac12\log0.64\approx0.2231\) natural units,
about \(0.3219\) bits. Converting both coordinates to millimetres adds
\(\log1000\) to each marginal entropy and \(2\log1000\) to the joint entropy,
so the difference is unchanged, as
\cref{thm:p15-w02b-p15-042-difference-invariance} requires.

\emph{Nonlinear coordinates.} In
\cref{thm:p15-w02b-p15-042-boundaries}, items 1--2, the uniform density on
\((0,1)\) is the maximum-entropy density for the \(x\)-coordinate by
\cref{thm:p15-w02b-p15-042-volume}, while in the \(y=x^2\) coordinate the
maximum is attained by the law with \(x\)-density \(2x\). A claim that one of
these two laws is ``more uncertain'' than the other is therefore not a claim
about the laws alone. Item~3 shows that a smooth monotone change of coordinate
can send a finite entropy to \(-\infty\).
\end{hbexample}

\begin{hbwarning}{def:p15-w02b-p15-042-limits}{What a continuous entropy does not measure}
\begin{enumerate}
\item The value of a continuous entropy, its sign, and the ranking of two laws
      by it depend on the reference measure and coordinates
      (\cref{thm:p15-w02b-p15-042-boundaries}). A maximum-entropy assignment
      on a bounded region is uniform only relative to the stated volume.
\item Its existence is not automatic: densities can have entropy \(+\infty\),
      \(-\infty\) or none, and laws without densities have none
      (\cref{thm:p15-w02b-p15-042-defined-cases}).
\item Several discrete facts fail: entropies and conditional entropies can be
      negative (\cref{thm:p15-w02b-p15-042-conditional}).
\item The coordinate-change formula needs the Jacobian identity and the
      integrability of \(p\log j\); without the latter the formula can
      produce an infinite value from a finite one
      (\cref{thm:p15-w02b-p15-042-boundaries}, item~3).
\item Entropy is a functional of a probability law. It is not the truth of a
      message, the evidential support for a hypothesis, or the value of a
      decision. Shannon sets the meaning of messages aside as irrelevant to
      the engineering problem \cite[p.~1]{Shannon1948MathematicalTheory}.
\end{enumerate}
\end{hbwarning}

\section{Relationships and sources}
\label{sec:p15-w02b-p15-042-relationships}

The bounded-interval maximum-entropy model with a fixed mean in
\cref{thm:p15-measurement-asymmetric-model} is the same Gibbs comparison as
\cref{thm:p15-w02b-p15-042-volume,thm:p15-w02b-p15-042-exponential}, with
an exponential comparison density on a bounded interval; its conclusion is
relative to the measurement coordinate in the same sense as item~8. Densities
relative to a general sigma-finite measure, their uniqueness and the change of
measure formula are those of \cref{thm:p15-radon-nikodym,thm:p15-density-calculus}.
With counting measure as reference, \(H_\mu\) is Shannon's discrete entropy,
for which the value is not relative to a coordinate system
\cite[sec.~20, p.~37]{Shannon1948MathematicalTheory}.

\emph{The continuous generalization in H\'ajek and Staffel.} Their chapter
states the finite principle of indifference as entropy maximization with
base-2 logarithms, says that entropy can be generalized to continuous
distributions, and states the generalized principle as choosing the
maximum-entropy law that meets a given constraint, naming Jaynes as an
advocate. Against both principles they set van Fraassen's factory of cubes
with side lengths up to one foot: spreading credence uniformly over side
length, over face area or over volume gives different answers to the
question whether a side is at most half a foot, and they judge that none of
these descriptions stands out as privileged
\cite[p.~244]{HajekStaffel2021Dynamics}; the handbook's account of that
passage is in \cref{sec:p15-rat41-synchronic}. In this chapter's terms, write
\(\ell\) for the side length in feet, so that the face area is \(\ell^2\)
and the volume \(\ell^3\). By \cref{thm:p15-w02b-p15-042-volume}, with no
constraint the maximum-entropy law relative to Lebesgue measure in the
coordinate \(\ell^b\) is the law uniform in \(\ell^b\), and
\cref{thm:p15-w02b-p15-042-power-coordinates} gives each of the three
uniform laws its entropy \(\log r+1-r\), with \(r=b/a\), in each of the three
coordinates. The laws uniform in \(\ell\), \(\ell^2\) and \(\ell^3\) give
the event \(\ell\leq\tfrac12\) the probabilities \(1/2\), \(1/4\) and \(1/8\),
and their entropies, in natural units and rounded to four places (divide by
\(\log2\approx0.6931\) for bits), are:
\begin{center}
\small\setlength{\tabcolsep}{3pt}
\begin{tabular}{lccc}
\hline
Law uniform in & \(H(\ell)\) & \(H(\ell^2)\) & \(H(\ell^3)\)\\
\hline
side \(\ell\) & \(0\) & \(\log2-1\approx-0.3069\) & \(\log3-2\approx-0.9014\)\\
area \(\ell^2\) & \(\tfrac12-\log2\approx-0.1931\) & \(0\) & \(\log\tfrac32-\tfrac12\approx-0.0945\)\\
volume \(\ell^3\) & \(\tfrac23-\log3\approx-0.4319\) & \(\tfrac13-\log\tfrac32\approx-0.0721\) & \(0\)\\
\hline
\end{tabular}
\end{center}
Each row is a single law written in three coordinates, and along a row the
entropy changes by the expected logarithm of the Jacobian factor, as in
\cref{thm:p15-w02b-p15-042-coordinate-change-thm}. That is the transport of
one fixed prior, the operation \cref{thm:p15-bayes-transformation} separates
from reimposing uniformity in the new coordinate. Down each column a
different law attains the maximum \(0\). The table therefore makes the
source's conflict exact without resolving it. The continuous maximum-entropy
principle inherits the choice of reference measure that Shannon's item~8
exposes, so it does not supply the privileged description whose absence the
source records. The table does not decide the further question the source
leaves open, whether that absence is fatal to the principle; it names
Norton, White and Novack as defenders \cite[p.~244]{HajekStaffel2021Dynamics}.
The source's informal gloss, that a uniform distribution is maximally
uncertain, holds for a finite outcome set
(\cref{thm:p15-rat41-indifference}); for continuous laws each row of the
table is maximally uncertain only in its own column. The same factory with a
declared mean, and the effect of a change of reference measure on the
maximizer, are treated in \cref{thm:p15-w02b-p15-044-cubes,thm:p15-w02b-p15-044-reference}.

\emph{Source.} C.~E. Shannon, \emph{A Mathematical Theory of Communication},
the corrected reprint of the 1948 Bell System Technical Journal papers:
logarithmic units pp.~1--2; Section~6 and the label convention for \(H(x)\)
pp.~10--11; band-limited ensembles as \(n\)-dimensional distributions
Section~19, pp.~34--35; Section~20, items 1--9 and the remarks on coordinate
standards, differences and negative values, pp.~35--38
\cite{Shannon1948MathematicalTheory}. The reformulation relative to a reference
measure, the integrability hypotheses, the Gibbs proofs replacing the
variational derivations, the explicit Jacobian hypothesis and the worked
boundaries are this chapter's.
